\documentclass{article}
\usepackage{amsmath, amssymb}
\usepackage[margin=1in]{geometry}

\title{Discovering Context Effects from Raw Choice Data: Notes}
\author{}
\date{}

\begin{document}
\maketitle
\section{Motivation}
 Let us try to imagine that someone has to get from IIT Madras to
the airport, and they have $4$ real choices like Car, Bike, taxi, bus..
And so, based on the cost, comfort, speed they have some natural
preference ranking among this $4$, let's say they choose or pick
\begin{equation*}
    \text{Car} > \text{Taxi} > \text{Bike} > \text{Bus}
\end{equation*}

Now, suppose we add a $5^{\text{th}}$, a clear irrelevant option to the
list which would almost take $3$ hours and nobody would seriously
choose this. Common sense tells us that adding `walking' doesn't
really change the feel about the car v/s the taxi v/s the bike, as it
is the irrelevant alternative, and their \textbf{relative preference} between all these $4$ alternative would not change at all.

So this idea is called independence of irrelevant alternatives and the
core behind the classical model of choice.

But the core idea is that maybe human brains may not be registering
the preference in the same way of that of multinomial logits. As
there is a strong experimental evidence proves that it would shift the attractiveness of things in the way that how attractive the existing options really looks $\rightarrow$ this is actually called violations of independence of irrelevant alternatives.\cite{seshadri2019}

\section{Independence of Irrelevant Alternatives}

$\mathcal{X}$ being the universal set, and $A$ and $B$ being the item
sets which are subsets of the universal set and The universal set X would have every choices fixed throught as well( in the above example the walking option as well because all the decision drawn for the real choices would come from this same fixed universe.). Also, let us say that $x, y$ are the
items of that set, then

\begin{equation}
    \frac{P(x \mid A)}{P(y \mid A)} = \frac{P(x \mid B)}{P(y \mid B)}
\end{equation}

\subsection*{My worked example}

So, now, let us call Car, Taxi, Bike, Bus, Walking as
$C, T, Bi, Bu, W$ for our own convenience; now let us work out
the example as well.

So,

Item \quad Symbol \quad Weight ($w$)

Car \quad\quad $C$ \quad\quad $0.40$

Taxi \quad\quad $T$ \quad\quad $0.25$

Bike \quad\quad $Bi$ \quad\quad $0.15$

Bus \quad\quad $Bu$ \quad\quad $0.12$

Walking \quad $W$ \quad\quad $0.08$

\hspace{1.6cm} Total \quad $1.00$

And the item set to work with:
\begin{align*}
    A &= \{C, T, Bi\}, \quad \sum w_A = 0.80 \\
    B &= \{C, T, W\}, \quad \sum w_B = 0.73
\end{align*}

to check whether

\begin{equation*}
    \frac{P(C \mid A)}{P(T \mid A)} = \frac{P(C \mid B)}{P(T \mid B)}
\end{equation*}

\[
P(C \mid A) = \frac{w_C}{\sum w_A} = \frac{0.40}{0.80}
\qquad\qquad
P(T \mid A) = \frac{w_T}{\sum w_A} = \frac{0.25}{0.80}
\]

\[
P(C \mid B) = \frac{w_C}{\sum w_B} = \frac{0.40}{0.73}
\qquad\qquad
P(T \mid B) = \frac{w_T}{\sum w_B} = \frac{0.25}{0.73}
\]

So,

\begin{equation*}
    \frac{P(C \mid A)}{P(T \mid A)} = \frac{P(C \mid B)}{P(T \mid B)}
\end{equation*}

or, \quad
$\dfrac{\left(\dfrac{0.40}{0.80}\right)}{\left(\dfrac{0.25}{0.80}\right)}
= \dfrac{\left(\dfrac{0.40}{0.73}\right)}{\left(\dfrac{0.25}{0.73}\right)}$

or, \quad
$\dfrac{0.40}{0.80} \times \dfrac{0.80}{0.25}
= \dfrac{0.40}{0.73} \times \dfrac{0.73}{0.25}$

or, \quad $\dfrac{0.40}{0.25} = \dfrac{0.40}{0.25}$

or, \quad $1 = 1$.

Hence proved.
\section{Multinomial Logits}

Mathematically,
\begin{equation*}
    P(x \mid C) = \frac{\exp(v(x))}{\displaystyle\sum_{y \in C} \exp(v(y))}
\end{equation*}

with the constraint
\begin{equation*}
    \sum_{y} v(y) = 0 \quad \rightarrow \quad \text{for the purpose of identifiability.(Shift invariant)}
\end{equation*}

So, let us also know the logic behind the purpose of identifiability with a example..

Suppose we have,
\begin{align*}
    v(C) &= 2.5 \\
    v(T) &= -0.9
\end{align*}

So,
\begin{align*}
    \exp(v(C)) &= \exp(2.5) = 12.18 \\
    \exp(v(T)) &= \exp(-0.9) = 0.407
\end{align*}

So,
\begin{equation*}
    P(C \mid \{C,T\}) = \frac{12.18}{12.18 + 0.407} = 0.968
\end{equation*}

While checking the constraint,
\begin{equation*}
    v(C) + v(T) = 1.6 \neq 0.
\end{equation*}

So, solving for $k$:
\begin{align*}
    (2.5+k) + (-0.9+k) &= 1.6 + 2k = 0 \\
    \text{or,} \quad k &= -0.8
\end{align*}

And after applying the shift,
\begin{align*}
    v(C) &= 1.7 \\
    v(T) &= -1.7
\end{align*}

Now, to verify that the probability is unchanged,
\begin{align*}
    P(C \mid \{C,T\}) &= \frac{\exp(1.7)}{\exp(1.7) + \exp(-1.7)} \\
    &= \frac{5.474}{5.657} \\
    &= 0.968.
\end{align*}

Hence proved. and moreover this method as in this constraint is used to know about the true one representative as 
2.5, -0.9 and 1.7, -1.7, the second one seems as better representation and is giving more idea about the uniquely defined behaviour.

\section{Universal logit model}
So here in this, instead of just a single utility item, it is defining and assigning different seperate utilities to each items. Mathematically, 
\begin{equation*}
    P(x \mid C) = \frac{\exp(u(x \mid C))}{\displaystyle\sum_{y \in C} \exp(u(y \mid C))}
\end{equation*}

with the constraint
\begin{equation*}
    \sum_{y \in C} u(y \mid C) = 0 \quad \forall\, C \subseteq \mathcal{X}
\end{equation*}

\section{Lemma 1}

Lemma 1 says that for any universal logits, the context dependent
utility $u(x \mid C)$ can be written as the sum of terms indexed by
subset $B$ where $B$ is the subset of $C \setminus \{x\}$. \cite{seshadri2019}

Mathematically,
\begin{equation*}
    u(x \mid C) = \sum_{B \subseteq C \setminus \{x\}} v(x \mid B)
    = v(x) + \sum_{y \in C \setminus x} v(x \mid \{y\})
    + \sum_{\{y, z\} \subseteq C \setminus x} v(x \mid \{y, z\}) + \dots
\end{equation*}

with the constraint:
\begin{equation*}
    \sum_{x \notin B} v(x \mid B) = 0 \quad \text{for every } B.
\end{equation*}

\subsection*{An example}

Set $A = \{C, T, Bi\}$.

Let us find out the subset for all first:
\begin{equation*}
    \{\emptyset\}, \; \{C\}, \; \{T\}, \; \{Bi\}, \; \{C, T\}, \;
    \{T, Bi\}, \; \{Bi, C\}, \; \{C, T, Bi\}.
\end{equation*}

Now:
\begin{align*}
    u(C \mid A) &= \underset{(1)}{v(C)} + \underset{(2)}{v(C \mid \{T\})} + \underset{(3)}{v(C \mid \{Bi\})} + \underset{(4)}{v(C \mid \{T, Bi\})} \\[1mm]
    u(T \mid A) &= \underset{(5)}{v(T)} + \underset{(6)}{v(T \mid \{C\})} + \underset{(7)}{v(T \mid \{Bi\})} + \underset{(8)}{v(T \mid \{C, Bi\})} \\[1mm]
    u(Bi \mid A) &= \underset{(9)}{v(Bi)} + \underset{(10)}{v(Bi \mid \{C\})} + \underset{(11)}{v(Bi \mid \{T\})} + \underset{(12)}{v(Bi \mid \{C, T\})}
\end{align*}

So, here the seven constraint equations:
\begin{align}
    v(C) + v(T) + v(Bi) &= 0 \\
    v(T \mid \{C\}) + v(Bi \mid \{C\}) &= 0 \\
    v(C \mid \{T\}) + v(Bi \mid \{T\}) &= 0 \\
    v(C \mid \{Bi\}) + v(T \mid \{Bi\}) &= 0 \\
    v(C \mid \{T, Bi\}) &= 0 \\
    v(T \mid \{C, Bi\}) &= 0 \\
    v(Bi \mid \{C, T\}) &= 0
\end{align}

Constraints (1)--(4) tie pairs of parameters together (4 constraints),
while constraints (5)--(7) directly kill single terms.

So, now the new equations after cutting these three terms would be:
\begin{align*}
    u(C \mid A) &= \underset{(1)}{v(C)} + \underset{(2)}{v(C \mid \{T\})} + \underset{(3)}{v(C \mid \{Bi\})} \\
    u(T \mid A) &= \underset{(4)}{v(T)} + \underset{(5)}{v(T \mid \{C\})} + \underset{(6)}{v(T \mid \{Bi\})} \\
    u(Bi \mid A) &= \underset{(7)}{v(Bi)} + \underset{(8)}{v(Bi \mid \{C\})} + \underset{(9)}{v(Bi \mid \{T\})}
\end{align*}

$\rightarrow$ 9 parameters.


\fbox{$9 - 4 = 5$ free parameters.}
% \section*{Maximum Likelihood Estimator}
% Before estimating the parameters of the codext dependent model via maximum likelihood estimation, I would like to do for the multinomial logit model first..

\section*{Maximum Likelihood Estimation: The Bradley Terry Model}

To find and understand the underlined mechanics of the model in the choice set, we will observe some data where we need to  assume a probabilistic model, and for now we are using Bradley terry model.
Actually, the goal of doing so is to find that unknown parameter i.e. the compressed representation of the observed data.
\section{Assumption followed for Estimation}
The Assumptions that one typically works with when we are dealing with the data about the observations are following :
\item Observations are independent 
mathematically, 
p(xi/xj)= P(xi) for all i not equals j.
\item Identically distributed.
that means, 
P(xi=1)=P(xj=1)=p for all i,j
\section{Likelihood Function}
In order to get a principled way to get estimates from the data, we need a likelihood function.
Likelihood function will give us the value of the parameter where the chance of seeing the data is highest.
As per Fisher's principle of maximum likelihood \citep{fisher1922mathematical}, given $n$
observations $x_1, x_2, \dots, x_n$ drawn from a distribution parameterised by $p$, the
likelihood function is defined as the joint probability of observing the data, viewed as a
function of the parameter:
\begin{align}
    L(p ; \{x_1, x_2, \dots, x_n\}) &= P(x_1, x_2, \dots, x_n ; p) \label{eq:likelihood-def} \\
    &= P(x_1 ; p) \cdot P(x_2 ; p) \cdots P(x_n ; p) \qquad \text{(by independence)} \label{eq:likelihood-product} \\
    &= \prod_{i=1}^{n} p^{x_i} (1-p)^{1-x_i} \label{eq:likelihood-bernoulli}
\end{align}
where the last line follows from taking each $x_i$ to be a Bernoulli trial with success
probability $p$, so that $P(x_i ; p) = p^{x_i}(1-p)^{1-x_i}$. The maximum likelihood estimate
$\hat{p}$ is then the value of $p$ that maximises Equation~\eqref{eq:likelihood-bernoulli}.

% \subsection*{Setting up the likelihood}

% Let us consider that we observe $N$ independent pairwise comparisons, where in the
% $n^{\text{th}}$ comparison item $i_n$ is chosen over item $j_n$. Under
% the multinomial logit model with $|C| = 2$, the probability of the
% observed winner is

% \begin{equation*}
%     P(i_n \succ j_n)
%     = \frac{\exp(v_{i_n})}{\exp(v_{i_n}) + \exp(v_{j_n})}
%     = \sigma(v_{i_n} - v_{j_n})
% \end{equation*}

% where $\sigma(x) = 1 / (1 + e^{-x})$ is the logistic function. So the
% likelihood of the full data is just the product over all $N$ comparisons

% \begin{equation*}
%     \mathcal{L}(\mathbf{v})
%     = \prod_{n=1}^{N} \sigma(v_{i_n} - v_{j_n}).
% \end{equation*}

% Taking log for the numerical convenience,

% \begin{equation*}
%     \ell(\mathbf{v})
%     = \log \mathcal{L}(\mathbf{v})
%     = \sum_{n=1}^{N} \log \sigma(v_{i_n} - v_{j_n}).
% \end{equation*}

% \subsection*{Gradient of the log-likelihood}

% Now, to find the MLE we would want to maximize $\ell(\mathbf{v})$, so let
% us take the partial derivative with respect to $v_k$.

% For a single comparison $(i, j)$,

% \begin{align*}
%     \frac{\partial}{\partial v_i} \log \sigma(v_i - v_j)
%         &= 1 - \sigma(v_i - v_j) \\
%     \frac{\partial}{\partial v_j} \log \sigma(v_i - v_j)
%         &= -\bigl(1 - \sigma(v_i - v_j)\bigr)
% \end{align*}

% and summing over all $N$ comparisons,

% \begin{equation*}
%     \frac{\partial \ell}{\partial v_k}
%     = \underbrace{W_k}_{\text{wins of } k}
%     \;-\;
%     \underbrace{\sum_{n \,:\, k \in \{i_n, j_n\}} P(k \text{ beats opponent})}
%         _{\text{expected wins of } k}
% \end{equation*}

% or, in short form,

% \begin{equation*}
%     \frac{\partial \ell}{\partial v_k} = W_k - E_k.
% \end{equation*}

% This gives a very clean interpretation $\rightarrow$ at the MLE, the
% observed wins equal the expected wins for every item. That is called
% the moment matching condition.

% \subsection*{Gradient descent update}

% Since $\ell$ is strictly concave in $\mathbf{v}$, gradient ascent will
% converge to the unique global maximum. So the update rule becomes

% \begin{equation*}
%     v_k^{(t+1)} = v_k^{(t)} + \eta \, \bigl(W_k - E_k^{(t)}\bigr)
% \end{equation*}

% where $\eta$ is the learning rate. And after every step we also enforce
% the identifiability constraint $\sum_k v_k = 0$ by subtracting the mean
% so that the parameters are uniquely defined (same shift trick as in the
% MNL section above).

% \subsection*{Numerical check on the same commute example}

% Using the weights from Section 2, that is
% $w = (0.40,\, 0.25,\, 0.15,\, 0.12,\, 0.08)$ for
% $(C, T, Bi, Bu, W)$, the corresponding centered true parameters are

% \begin{equation*}
%     \mathbf{v}_{\text{true}}
%     = [\,0.853,\ 0.383,\ -0.128,\ -0.351,\ -0.757\,].
% \end{equation*}

% We simulate $N = 2000$ pairwise comparisons: for each comparison, a pair
% $(i, j)$ is drawn uniformly at random from the $5$ items, and the winner
% is drawn from Bernoulli$\bigl(\sigma(v_i - v_j)\bigr)$.

% Fitting the model in two independent ways:

% \begin{itemize}
%     \item {I-LSR} via \texttt{choix.ilsr\_pairwise} $\rightarrow$
%     Maystre's spectral method (Chapter 3 of his thesis).
%     \item {Gradient descent from scratch} $\rightarrow$ the update
%     rule above, run for $1000$ iterations with $\eta = 0.01$ on the
%     normalized gradient.
% \end{itemize}

% gives the estimates

% \begin{align*}
%     \mathbf{v}_{\text{I-LSR}} &= [\,0.827,\ 0.461,\ -0.122,\ -0.418,\ -0.749\,] \\
%     \mathbf{v}_{\text{GD}}    &\approx \mathbf{v}_{\text{I-LSR}}
% \end{align*}

% with mean absolute error $\approx 0.04$ against the true parameters.
% So, both the methods agree with each other and with the truth up to
% sampling noise. This is expected because the Bradley-Terry log-likelihood
% is strictly concave $\rightarrow$ any correct optimizer must land at the
% same unique MLE.

% Hence proved.





\bibliographystyle{plain}
\bibliography{references}

\end{document}